\abstract{ARC-ADHD defines recovery as a measurable return to stable goal-directed performance after a controlled challenge and regulation sequence. This companion study proposes synchronized wearable electroencephalography (EEG) with Muse S Athena as an observational measurement layer during ARC-ADHD use. The primary objective is to determine whether event-locked EEG spectral dynamics provide reproducible within-child information about challenge reactivity and post-regulation recovery beyond behavioral and contextual telemetry. The initial feature family includes prespecified spectral power in conventional delta, theta, alpha, beta, and low-gamma ranges, with primary emphasis on theta, alpha, and beta and explicit artifact controls. Recovery-related EEG features are normalized to a target-rule-matched calibration reference rather than the immediate preswitch state, preventing response-mapping changes from being misinterpreted as neural recovery. A candidate neural recovery time is defined from recovery-monitoring onset and compared with the independently constructed behavioral recovery latency. The study is explicitly not an EEG diagnostic study, a neurofeedback intervention, or evidence of clinical efficacy. Progression to EEG-informed adaptation is justified only if synchronization, signal quality, within-child reproducibility, and participant-held-out incremental value are demonstrated.}
